%%
%% This is file `sample-ccs2026.tex',
%% it has been extracted from `sample-sigconf.tex',
%% generated with the docstrip utility.
%%
%% It has been annotated with guidelines for submissions to CCS 2026
%% Most optional parts have been removed to provide a MINIMAL file
%% that needs to be used.
%% All specific CCS instructions are indicated with the keyword CCS
%%
%% For help and more latex instructions, refer to
%% `sample-sigconf.tex', provided in the distribution
%% https://portalparts.acm.org/hippo/latex_templates/acmart-primary.zip 
%%

\documentclass[sigconf]{acmart} %% CCS: DO NOT REMOVE
%%
%% \BibTeX command to typeset BibTeX logo in the docs
\AtBeginDocument{%
  \providecommand\BibTeX{{%
    Bib\TeX}}}

%% Rights management information.  This information is sent to you
%% when you complete the rights form.  CCS: These commands have SAMPLE
%% values in them; you MUST leave these commands in the CCS 2026
%% submission version. CCS: Upon acceptance, it is your responsibility as an author to replace
%% the commands and values with those provided to you when you
%% complete the rights form.
\setcopyright{acmlicensed} %% CCS: DO NOT REMOVE
\copyrightyear{2018} %% CCS: DO NOT REMOVE
\acmYear{2018} %% CCS: DO NOT REMOVE
\acmDOI{XXXXXXX.XXXXXXX} %% CCS: DO NOT REMOVE
\acmConference[Conference acronym 'XX]{Make sure to enter the correct
  conference title from your rights confirmation email}{June 03--05,
  2018}{Woodstock, NY}  %% CCS: DO NOT REMOVE
\acmISBN{978-1-4503-XXXX-X/2018/06}  %% CCS: DO NOT REMOVE

%%
%% end of the preamble, start of the body of the document source.
\begin{document}

%%
%% The "title" command has an optional parameter,
%% allowing the author to define a "short title" to be used in page
%% headers.

\title{The Name of the Title Is Hope} %% CCS: you MUST provide a title

%%
%% The "author" command and its associated commands are used to define
%% the authors and their affiliations.
%% Of note is the shared affiliation of the first two authors, and the
%% "authornote" and "authornotemark" commands
%% used to denote shared contribution to the research.

%% CCS: at submission time, the submission MUST be anonymized. Hence
%% authors MUST be commented out.

% \author{Ben Trovato}
% \authornote{Both authors contributed equally to this research.}
% \email{trovato@corporation.com}
% \orcid{1234-5678-9012}
% \author{G.K.M. Tobin}
% \authornotemark[1]
% \email{webmaster@marysville-ohio.com}
% \affiliation{%
%   \institution{Institute for Clarity in Documentation}
%   \city{Dublin}
%   \state{Ohio}
%   \country{USA}
% }

% \author{Lars Th{\o}rv{\"a}ld}
% \affiliation{%
%   \institution{The Th{\o}rv{\"a}ld Group}
%   \city{Hekla}
%   \country{Iceland}}
% \email{larst@affiliation.org}

% \author{Valerie B\'eranger}
% \affiliation{%
%   \institution{Inria Paris-Rocquencourt}
%   \city{Rocquencourt}
%   \country{France}
% }

% \author{Aparna Patel}
% \affiliation{%
%  \institution{Rajiv Gandhi University}
%  \city{Doimukh}
%  \state{Arunachal Pradesh}
%  \country{India}}

% \author{Huifen Chan}
% \affiliation{%
%   \institution{Tsinghua University}
%   \city{Haidian Qu}
%   \state{Beijing Shi}
%   \country{China}}

% \author{Charles Palmer}
% \affiliation{%
%   \institution{Palmer Research Laboratories}
%   \city{San Antonio}
%   \state{Texas}
%   \country{USA}}
% \email{cpalmer@prl.com}

% \author{John Smith}
% \affiliation{%
%   \institution{The Th{\o}rv{\"a}ld Group}
%   \city{Hekla}
%   \country{Iceland}}
% \email{jsmith@affiliation.org}

% \author{Julius P. Kumquat}
% \affiliation{%
%   \institution{The Kumquat Consortium}
%   \city{New York}
%   \country{USA}}
% \email{jpkumquat@consortium.net}

%%
%% By default, the full list of authors will be used in the page
%% headers. Often, this list is too long, and will overlap
%% other information printed in the page headers. This command allows
%% the author to define a more concise list
%% of authors' names for this purpose.
\renewcommand{\shortauthors}{Trovato et al.}

%%
%% The abstract is a short summary of the work to be presented in the
%% article.
\begin{abstract} %% CCS: an abstract MUST be provided.
  %% CCS: REMOVE the content below and replace it with appropriate content
  This document exemplifies the minimal information that must be kept
  for CCS 2026 submissions. The comments in the .tex file clarify
  which parts must be kept as such and which parts may/should be updated.
\end{abstract}

%%
%% The code below is generated by the tool at http://dl.acm.org/ccs.cfm.
%% Please copy and paste the code instead of the example below.
%%
\begin{CCSXML} %% CCS: DO NOT REMOVE this part. You MAY update the
               %% concepts.
<ccs2012>
 <concept>
  <concept_id>00000000.0000000.0000000</concept_id>
  <concept_desc>Do Not Use This Code, Generate the Correct Terms for Your Paper</concept_desc>
  <concept_significance>500</concept_significance>
 </concept>
 % <concept>
 %  <concept_id>00000000.00000000.00000000</concept_id>
 %  <concept_desc>Do Not Use This Code, Generate the Correct Terms for Your Paper</concept_desc>
 %  <concept_significance>300</concept_significance>
 % </concept>
 % <concept>
 %  <concept_id>00000000.00000000.00000000</concept_id>
 %  <concept_desc>Do Not Use This Code, Generate the Correct Terms for Your Paper</concept_desc>
 %  <concept_significance>100</concept_significance>
 % </concept>
 % <concept>
 %  <concept_id>00000000.00000000.00000000</concept_id>
 %  <concept_desc>Do Not Use This Code, Generate the Correct Terms for Your Paper</concept_desc>
 %  <concept_significance>100</concept_significance>
 % </concept>
</ccs2012>
\end{CCSXML}

\ccsdesc[500]{Do Not Use This Code~Generate the Correct Terms for Your Paper}
% \ccsdesc[300]{Do Not Use This Code~Generate the Correct Terms for Your Paper}
% \ccsdesc{Do Not Use This Code~Generate the Correct Terms for Your Paper}
% \ccsdesc[100]{Do Not Use This Code~Generate the Correct Terms for Your Paper}

%%
%% Keywords. The author(s) should pick words that accurately describe
%% the work being presented. Separate the keywords with commas.
\keywords{Do, Not, Use, This, Code, Put, the, Correct, Terms, for,
  Your, Paper} %% CCS: DO NOT REMOVE but you MAY update

% \received{20 February 2007} 
% \received[revised]{12 March 2009}
% \received[accepted]{5 June 2009}

%%
%% This command processes the author and affiliation and title
%% information and builds the first part of the formatted document.
\maketitle

\section{Introduction} %% CCS: You MAY change the title and,
                       %% obviously, add text, sections, figures,
                       %% tables, etc. 
%% CCS: For help and more latex examples, refer to
%% `sample-sigconf.tex', provided in the distribution
%% https://portalparts.acm.org/hippo/latex_templates/acmart-primary.zip 
%%

%%
%% The acknowledgments section is defined using the "acks" environment
%% (and NOT an unnumbered section). This ensures the proper
%% identification of the section in the article metadata, and the
%% consistent spelling of the heading.

%% CCS: to preserve anonymity, NO acknowledgements to fundings, projects or persons should be used at
%% submission time
%% CCS: this section MAY be used to acknowledge the use of AI when used only for minor editorial improvements (e.g., grammar, spelling, or light style polishing) 
% \begin{acks}
% This paper was edited for grammar using [Tool Name].
% \end{acks}

%%
%% The next two lines define the bibliography style to be used, and
%% the bibliography file.
\bibliographystyle{ACM-Reference-Format}
\bibliography{sample-base}


%%
%% Appendices
\appendix %% CCS: DO NOT REMOVE

\section{Open Science} %% CCS: DO NOT REMOVE

%% CCS: REMOVE the content below and replace it with appropriate content
Each submitted paper must include an “Open Science” appendix that:
\begin{itemize}
\item Enumerates all artifacts needed to evaluate the paper’s core contributions (e.g., code, datasets, models, configuration files, scripts, documentation, benchmarks).
\item Clearly describes how the program committee can access each artifact during double-blind review (including anonymous URLs or credentials, where applicable).
\item Explicitly justifies any artifact that cannot be shared (e.g., due to licensing restrictions, responsible disclosure concerns, safety or privacy of study subjects, or deployment risks if adversarial methods are released prematurely). When full sharing is not possible, authors are encouraged to provide partial, synthetic, or redacted artifacts that still allow reviewers to assess the methodology.
\item In case no artifact is needed to evaluate the paper’s core contributions, the authors should state it explicitly.
\end{itemize}

\section{Ethical Considerations} %% CCS: the entire sectin MAY BE REMOVED only if it is
                                %% clear that your
                       %% paper does not raise any ethical
                       %% concerns. In case of doubts, keep the
                       %% section and explains why your paper  does
                       %% not raise any ethical concerns. 

 %% CCS: REMOVE the content below and replace it with appropriate content
Authors are expected to consider the ethical implications and potential societal impact of their work. Papers that raise ethical concerns, such as those involving human subjects, user data, or real-world vulnerability analysis, must include a dedicated "Ethical Considerations" section. This section should discuss the balance of risks vs. benefits and the steps taken to minimize potential harm (e.g., responsible disclosure, data anonymization). Note that institutional (IRB/ERB) approval is neither strictly necessary nor always sufficient to demonstrate ethical conduct; we expect authors to reason about the ethics of their work beyond ensuring institutional compliance. For detailed guidance on community standards, we follow the USENIX Security'26 Ethics Policy\footnote{\url{https://www.usenix.org/conference/usenixsecurity26/call-for-papers\#ethics}}. This section does not count toward the page limit and must be placed after the 12-page main content.

\end{document}
\endinput
%%
%% End of file `sample-sigconf.tex'.
